% ======================================================================
% 第 I 节：绪论
% ======================================================================
\section{绪论}\label{sec:intro}

夸克组元图像（quark-constituent picture）在共振态和深度非弹性电子、中微子过程
两方面所取得的成功，使我们很难相信夸克并不存在。问题在于夸克从未被观测到。
这表明夸克出于某种原因不能作为独立粒子出现在末态中。关于这一点人们已经提出了
许多推测。${}^{1}$

与夸克问题无关，施温格（Schwinger）多年前便观察到：${}^{2}$ 如果真空极化完全
屏蔽了规范理论中的电荷，那么规范理论的矢量介子可以获得非零质量。施温格用
一空间一时间的量子电动力学（QED）的精确解阐明了这一结果，其中光子对于任何
非零电荷 $e$ 获得质量 $\sim e^{2}$ [在这一理论中 $e$ 具有 $(\mathrm{mass})^{1/2}$
的量纲]。施温格猜测相同的效应在四维中对于足够强的耦合也可能发生。

洛文斯坦（Lowenstein）和斯维卡（Swieca）${}^{3}$ 以及卡谢尔（Casher）、
科古特（Kogut）和苏斯坎德（Susskind）${}^{4}$ 对施温格模型的进一步研究表明，
该模型的渐近态只包含有质量的光子，而不含电子。尽管如此，正如 Casher 等人
所详细证明的，电子出现在深度非弹性过程中，并且在短时间、短距离上表现得像
自由点状粒子。阻止电子出现在末态中的极化效应发生在更长的时间尺度上
（长于 $1/m_{\gamma}$，其中 $m_{\gamma}$ 是光子质量）。

本文将提出一种使夸克保持束缚的新机制。该机制仅适用于规范理论。我们将在四维
时空的规范理论的强耦合极限下说明这一机制。然而，这里讨论的模型带有内禀的
紫外截断，并且在强耦合极限下所有粒子质量（包括规范场质量）都远大于截断；
其结果是理论远远偏离协变性。

这里提出的禁闭机制是软的（长时标）。然而，在本文讨论的模型中，截断破坏了夸克
具有自由点状行为的可能性。

本文所讨论的模型是在四维欧几里得格点上建立起来的规范理论。格点间距 $a$ 的
倒数充当紫外截断。采用欧几里得空间（即虚时间而非实时间）代替洛伦兹空间并非
严重限制；格点理论的能量本征态（包括散射态）可以通过``转移矩阵''（transfer-matrix）
形式化来确定，这一点已由 Suri${}^{5}$ 讨论并由 Wilson 和 Kogut${}^{6}$ 综述。
第三节的第 B 小节将简要讨论转移矩阵方法。

在施温格关于四维的猜想中，对于任何小于临界耦合 $e_{c}$ 的电荷 $e$，光子质量都为零；
当 $e > e_{c}$ 时，光子质量非零并随 $e$ 变化。图~\ref{fig:1} 展示了 $m_{\gamma}$
对 $e$ 的曲线可能的样子。$e_{c}$ 是一个非解析点。类似的非解析点称为临界点，
出现在凝聚态物理中某些类型的相变处。例如，考虑无外场时的铁磁体。对于任何
高于居里温度 $T_{c}$ 的温度，自发磁化强度 $M$ 为零。在 $T_{c}$ 以下，$M$ 非零
并且是温度的函数。在 $T_{c}$ 处可能发生一阶相变（此时 $M$ 在 $T_{c}$ 处不连续），
也可能发生二阶相变（临界点），此时 $M \to 0$ 当 $T \to T_{c}$（从两侧趋近）。

\begin{figure}[t]
  \centering
  \includegraphics[width=0.42\textwidth]{fig1}
  \refstepcounter{figure}%
  \caption*{图 1. 光子质量对重正化电荷 $e$ 的推测性曲线，单位未知。
    $e_{c}$ 处的相变是二阶的（见正文）。}
  \label{fig:1}
\end{figure}

通过与凝聚态情形的类比，可以把光子质量从零到非零的转变看作是一种相变：
这一类比最好这样理解——把量子电动力学的粒子想象为一种介质（以太）的激发。
在这种情况下，正是以太在 $e_{c}$ 处发生相变。这里再次出现一个问题：这一相变
是一阶的（参看图~\ref{fig:2}）还是二阶的（图~\ref{fig:1}）。（Coleman 和
Weinberg${}^{7}$ 在另一背景下发现了一阶相变的非平凡例子。）

本文讨论的模型是一个阿贝尔规范场与有质量夸克耦合（强度为 $g$）。在弱耦合下，
规范场表现得像一个普通的自由零质量场（尽管格点量子化引入了修正），夸克不束缚。
在强耦合下规范场有质量，夸克被束缚，表明存在第二个相。因此在 $g$ 的某个中间值
处应当存在相变。目前对这一相变一无所知。

\begin{figure}[t]
  \centering
  \includegraphics[width=0.42\textwidth]{fig2}
  \refstepcounter{figure}%
  \caption*{图 2. 若在 $e_{c}$ 处发生一阶相变时，光子质量对重正化电荷 $e$ 的
    推测性曲线。}
  \label{fig:2}
\end{figure}

本文描述的量子化程序和强耦合近似同样适用于非阿贝尔规范理论。这一点将在第三节
中简要说明。

格点理论的强耦合展开（见第四节）的一个非凡特征是，它具有与强子相对论性弦模型
相同的总体结构。${}^{8}$ 规范理论的真空期望值（在强耦合展开中）涉及对所有夸克
路径以及连接这些路径的所有曲面求和；这些曲面由强耦合处理下的规范场产生。路径
和曲面定义在离散格点上。在把格点上的曲面与弦模型的连续曲面联系起来时存在几何
上的困难；目前尚不知道这些困难能否被克服。

第二节将定性讨论夸克束缚机制的性质。第三节将在离散格点上经典地、量子力学地
表述规范理论。第四节解释格点规范理论的强耦合展开。第五节对弱耦合作粗略讨论。
第六节简要讨论洛伦兹不变性问题以及与弦模型的关系。
